\documentclass[10pt]{article}

\usepackage[margin=1in]{geometry}
\usepackage[T1]{fontenc}
\usepackage{amsmath,amssymb}
\usepackage{graphicx}
\usepackage{booktabs}
\usepackage{float}
\usepackage{xcolor}
\usepackage[colorlinks=true,allcolors=blue]{hyperref}
\IfFileExists{natbib.sty}{\usepackage[numbers,sort&compress]{natbib}}{%
  \providecommand{\citep}[1]{[#1]}\providecommand{\citet}[1]{[#1]}%
  \providecommand{\cite}[1]{[#1]}}

% A placeholder that is impossible to overlook in a draft, and impossible to typeset by accident
% in a submission: every one of these has a line in PLACEHOLDERS.md saying which experiment
% supplies it and what we expect it to show.
\newcommand{\TBD}[2]{\textcolor{red}{\textbf{[TBD-#1: #2]}}}
% A number that is measured. Used so a reader (and a co-author) can see at a glance which figures
% in a sentence are evidence and which are prose.
\newcommand{\m}[1]{#1}

\title{\textbf{Design Intent as an Evaluable Output:\\
Pocket-Grounded Hit-to-Lead Optimisation with Language Models}}
\author{}
\date{}

\begin{document}
\maketitle

\begin{abstract}
Medicinal chemists optimise a hit by acting on an opportunity in the binding pocket, and every
analogue they make encodes a design hypothesis: an unsatisfied donor to reach, an empty lipophilic
shelf to fill, a clash to relieve. We ask whether a language model can be taught that hypothesis
and then act on it. We factorise hit-to-lead design into a \emph{reason} model $R$ (complex
$\rightarrow$ design intent, in words) and a \emph{molecule} model $M$ (complex $+$ intent
$\rightarrow$ analogue), each a separate adapter over a chemistry LLM, both conditioned on a
per-atom contact table for that specific ligand in that specific pocket. Training uses
crystal-derived analogue pairs with agent-written intents, within-assay ChEMBL pairs carrying
measured potency deltas, and design intents mined from the text of \m{1{,}622} published SAR papers
(\m{17{,}601} pairs; \m{39.3\%} of extracted rationales state an instruction rather than restate an
outcome). Finetuning lifts a chemistry LLM from \m{0.118} to \m{0.973} SMILES validity and from
\m{3.9\%} to \m{46.5\%} sound reasoning under a fixed model-labeller panel. On a molecule-level
holdout, the reference intent recovers the analogue actually made at \m{0.919} mean
Tanimoto (median \m{1.000}) against \m{0.541} under an intent belonging to a different pair
($+0.379$ paired, $p=5\times10^{-50}$) --- though those reference intents were written with the
analogue in view and name a group that actually changed in \m{86.1\%} of rows, so that arm bounds
transcription fidelity rather than design. The deployable $R\!\rightarrow\!M$ pipeline reaches
\m{0.570} --- it retains only \m{8.7\%} of that gap, so the reason model, not the molecule model,
is the bottleneck. Withholding the contact table collapses what the explanation may cite
(groundedness \m{1.000} to \m{0.432}, residue specificity \m{0.388} to \m{0.121}) and leaves the
molecule essentially unchanged (\m{0.679} vs.\ \m{0.680}); a deliberately wrong pocket costs
\m{0.013} Tanimoto. The pocket governs what the model may claim, not yet what it makes, and we
report that asymmetry as the central result rather than as a limitation. On a publication-date
retrospective $R\!\rightarrow\!M$ recovers \TBD{3}{retrospective recovery}, and a blinded chemist
panel rates $R$'s intents at \TBD{1}{chemist panel}.
\end{abstract}

\section{Introduction}

Hit-to-lead optimisation turns an initial binder into a lead through rounds of closely related
analogues. The enumeration is not the hard part: a matched-pair tool will hand a project team
thousands of analogues of any hit~\citep{wawer2014whatnext}. The hard part is choosing the five to
make this month and defending the choice, and chemists make that choice by reading the complex ---
a donor left unsatisfied, a subpocket left empty, a carbon pressed against a polar rim. Each
modification is a hypothesis about the pocket, and the measured outcome of the previous round is
what informs the next.

Two literatures address parts of this and meet nowhere in the middle. Pocket-conditioned generative
models place atoms inside a measured binding site and are, by construction, interaction-aware;
some predict contacts or pharmacophores explicitly~\citep{peng2022pocket2mol,guan2023targetdiff,
schneuing2024diffsbdd,feng2024lingo3dmol,jiang2024pocketflow}. None of them says \emph{why}. On the
other side, chemistry language models now propose and justify modifications in natural language,
increasingly with explicit reasoning traces trained from outcome
data~\citep{zhao2024chemdfm,narayanan2025ether0,cavanagh2024smileyllama,liu2026desro,liu2026drugr,
wen2026logos}, and recent work does put pockets and even required residue contacts into an LLM
prompt~\citep{macdougall2026spatial}. What is still missing is the thing a project team actually
consumes: a stated design hypothesis whose claims can be \emph{checked}, one by one, against the
structure the model was shown.

That check is the contribution of this paper. We condition on a per-atom contact table --- for this
ligand, in this pocket, which residue atom sits how far from which ligand atom, with interaction
typing --- and we treat the design intent as a first-class output that is evaluated separately from
the molecule. Because the intent is a sentence that names residues, every clause it asserts can be
matched against the rows of the table the model received, and a sentence that invents a contact can
be caught automatically and counted.

Making the intent evaluable also makes it \emph{replaceable}. We factorise
$p(y,z\mid x,c)=p_R(z\mid x,c)\,p_M(y\mid x,c,z)$, with $x$ the hit, $c$ its contact table, $z$ the
intent and $y$ the analogue, and train $R$ and $M$ as separate adapters with separate losses. The
factorisation is the chain rule, so it costs nothing in expressiveness; what it buys is separate
supervision and separately trained weights for the two factors. It does not buy the bottleneck
itself: a jointly trained model that emits the sentence before the SMILES is already routing the
molecule through the reason in token order, which is why an intent can be teacher-forced into it and
why the comparison between joint and factorised is an empirical question rather than an
architectural one. We
can then intervene on $z$ alone --- substituting the reference intent, a sentence written for
a different pair, or no sentence at all --- and read off how much of the molecule the sentence
decides. Prior reason-then-molecule systems generate the rationale and the structure inside one
model and one trace~\citep{liu2026desron,liu2026drugr,wen2026logos}, so the reason cannot be swapped
and this measurement is unavailable; earlier work factorised \emph{substructure}-to-molecule rather
than sentence-to-molecule~\citep{jin2020rationale}.

Our supervision is also different in kind. Reasoning traces in recent molecular-optimisation work
are manufactured from the outcome: derived by analysing which fragment groups improved a measured
or computed property~\citep{liu2026desro,liu2026drugr}. We instead mine the sentence the authors
themselves wrote --- the hypothesis as stated in the paper, before the assay result was known to the
reader --- from \m{1{,}622} medicinal-chemistry papers, paired with the two compounds the paper
labels and with their potency difference measured \emph{in the same assay} of that same paper. To
our knowledge no prior work extracts the stated design rationale of a medicinal-chemistry
modification from paper text; the existing literature-mining work extracts reactions, procedures and
entities~\citep{chemminer2024,materials2025ie}. We also report, in the same breath, that only
\m{39.3\%} of what we extract states an instruction rather than restating an outcome, which puts the
majority of our own corpus in the same category as the methods we are distinguishing ourselves from.

\paragraph{What we find, including what we did not want to find.}
\begin{enumerate}\itemsep2pt
\item \textbf{An intent sentence behaves like an executable specification.} Given the reference
intent, $M$ reproduces the analogue made at \m{0.919} mean Tanimoto, median \m{1.000}, on a
molecule-level holdout; against a sentence belonging to another pair, \m{0.541} ($+0.379$ paired,
$p=5\times10^{-50}$). A \emph{wrong} sentence is worse than no sentence (\m{0.147} below,
$p=5\times10^{-40}$) and costs validity (\m{100\%}$\rightarrow$\m{91.5\%}). The model follows the
instruction; it does not merely decorate.
\item \textbf{The no-intent arm is not a baseline.} With no sentence, $M$ returns the input
unchanged in \m{100\%} of rows. Every property delta is exactly zero and its Tanimoto to the
reference is the \emph{input's} own similarity. We report it as a diagnostic --- zero intent, zero
edit --- and use the mismatched arm as the honest control.
\item \textbf{The pocket governs the prose, not the molecule.} Withholding the contact table takes
groundedness from \m{1.000} to \m{0.432} and residue specificity from \m{0.388} to \m{0.121}, and
leaves similarity to the reference untouched (\m{0.679} vs.\ \m{0.680}). A deliberately wrong pocket
costs \m{0.013} Tanimoto at cohort scale. We therefore claim \emph{grounded attribution}, not
pocket-aware generation, and we identify the likely mechanism: the fields that would describe
\emph{where a group could go} --- open room, slot shape, exit vectors --- are computable but are
excluded from the model's input because, computed only at the edit site, they leaked it.
\item \textbf{The reason model is the bottleneck.} $R\!\rightarrow\!M$ reaches \m{0.570}, retaining
\m{8.7\%} of the oracle-over-mismatch gap, and $R$'s intents are repetitive (uniqueness \m{0.238}).
The factorisation is validated as an upper bound and as a product claim --- a chemist can type the
intent today --- not as an end-to-end gain.
\item \textbf{Property-directed intents move the property, unevenly.} Polarity ($+3.62$ TPSA,
$n=259$), lipophilicity ($-0.63$ cLogP, $n=146$) and size-down ($-31.2$\,Da, $n=75$) move as asked
and are flat or reversed under a mismatched intent; size-up is flat and rigidity moves the wrong
way.
\end{enumerate}

We also publish the corpus defects, the metrics we refuse to report, and two numbers from this
programme that were withdrawn after failing to reproduce. The cheapest credibility in a paper of
this kind is naming your own retractions.

\section{Methods}

\subsection{Task, and why the unit of structure is a complex}

The model receives a hit molecule $x$ and a structured description $c$ of that molecule's
environment in its target, and emits an intent sentence $z$ and an analogue $y$. Critically, $c$ is
\emph{ligand-specific}: one block per ligand heavy atom, sixteen fields giving the nearest protein
atom and its distance, a typed interaction with hydrogen-bond angle, a second polar partner, water
counts and nearest water, a lining description, an open fraction and a burial reading (full field
list in the SI). A target-level pocket is therefore \emph{not} sufficient input: knowing the protein
buys nothing for a ligand whose pose we have never computed. The unit of work is the pair (molecule, target) --- a
complex --- and the cost of the structural conditioning is one structure per molecule, not one per
target. This single fact drives the dataset economics reported below.

Where a crystal structure of the complex exists we use it. Where it does not --- which is the case
for essentially every molecule in the literature corpus --- we cofold the complex with
Boltz~\citep{wohlwend2024boltz1,passaro2025boltz2} at one diffusion sample and derive the table from
the returned structure. This bounds what may be claimed from it, and we state the bound rather than
hoping a referee misses it:

\begin{quote}
Where no crystal structure exists, the complex is a single-sample cofold. Its contact table is a
hypothesis about which residues line the ligand, not a measurement of distances. No distance taken
from a cofolded table is quoted as evidence in this paper, and covalent geometry is never read from
a cofold, because an unconstrained cofold forms no bonds and a warhead lands wherever the sampler
puts it.
\end{quote}

\paragraph{Covalent typing.} A covalent bond is asserted only where the table types the row
\texttt{COVALENT BOND}, requiring the modelled heavy-atom separation to lie within 0.35\,\AA\ of the
ideal length for that element pair (1.81\,\AA\ for C--S; the SARS-CoV-2 main protease complex
7JQ2 measures exactly 1.81\,\AA\ at CYS145 and is our calibration case). Separations of
2.16--2.71\,\AA\ are typed and disclosed as \emph{probable} and are never cited as a bond; anything
beyond 2.71\,\AA\ is a non-bonded contact however a sentence describes it. This matters because an
audit of an earlier generation of our own labels found that \m{26\textrm{--}40\%} of clauses calling
a residue a ``covalent anchor'' were in fact citing a 2.7--3.8\,\AA\ non-bonded contact --- a
failure of the text layer, not of the structure layer, which had typed those rows correctly all
along. Every clause matching \texttt{covalent|bond|anchor|warhead|adduct} is now gated on a
\texttt{COVALENT BOND} row in that pair's own table, and the gate runs irrespective of whether the
sentence quotes a distance.

\subsection{Corpora}

\begin{table}[H]\centering\small
\begin{tabular}{llll}
\toprule
corpus & what it is & scale & intent written by \\
\midrule
PocketPair-9k/10k/14k & PDB-derived analogue pairs with pockets & 7{,}692 / 9{,}453 / 12{,}392 rows & an agent, from the pair \\
PocketPair-21k & ChEMBL within-assay pairs, cofolded, with $\Delta$potency & 3{,}992 complexes & an agent, from the pair \\
PocketPair-Literature & design intents quoted from paper text & 1{,}622 papers, 17{,}601 pairs & the paper's authors \\
\bottomrule
\end{tabular}
\caption{The three supervision regimes. The distinction in the last column is the one that matters
for interpreting our strongest result, and we return to it in \S\ref{sec:leakage}.}
\end{table}

\paragraph{Literature extraction.} For each paper we extract design \emph{episodes}: a reference
compound, one or more analogues, what changed, the stated reason, a verbatim evidence quote, the
endpoint in the paper's own words, and the direction of the outcome. Molecules are not recovered by
substructure search --- which finds the authors' actual pair only \m{32\%} of the time --- but by
ChEMBL's own curation under the paper's compound labels. A potency $\Delta$ is computed only between
two molecules measured in the \emph{same assay} of that paper. Under the final brief (v5.3), on all
\m{1{,}622} papers, the endpoint is stated in \m{100\%} of records, the reference source in
\m{100\%}, a rationale in \m{100\%} and an evidence quote in \m{100\%}; extraction-meta leakage is
\m{0.00\%}; repeated boilerplate in \texttt{what\_changed} is \m{0.0\%}; and edits spanning multiple
sites fell from \m{44.8\%} to \m{5.5\%} between brief versions.

\paragraph{Corpus defects, stated here rather than in a footnote.} These are SAR pairs, not matched
molecular pairs: the gate is a similarity threshold plus a same-assay potency delta, and the
extracted transformation field is descriptive only. A 100-paper QA round found molecules contradicted
by the text in \m{3.3\%} of records and reasons honest in \m{90\%} --- but \m{60.7\%} of those
reasons \emph{restate an outcome} rather than state an instruction, so only \m{39.3\%} of the corpus
is the supervision we actually want. Training on outcome restatements teaches post-hoc narration,
which is precisely the failure an independent cohort round measured at explanation-sound \m{0.08}.
We therefore report the instruction subset separately \TBD{13}{instruction-subset training};
hypothesis: the instruction subset trains a more useful $R$ at equal token count. Separately,
explained edits carry \emph{small} potency deltas, so selecting pairs by $|\Delta|$ systematically
discards them --- intent mining has to precede pair selection, not follow it.

\paragraph{Complex coverage.} The \m{17{,}601} literature pairs require \m{6{,}590} distinct
complexes across \m{658} targets, of which \m{27} were already held. \m{552} targets have a
sequence in ChEMBL and \m{106} do not, placing \m{1{,}995} complexes and \m{4{,}967} pairs out of
reach without another sequence source. \m{369} complexes exceed our \m{1{,}187}-residue cofold cap
and are recorded in a separate file rather than dropped --- an earlier run lost an entire major
target family to exactly this cap, silently. The remaining \m{4{,}199} complexes are queued in
descending order of pairs unlocked, so a partial run is the most valuable slice rather than a random
one: the first \m{1{,}000} complexes unlock \m{43.2\%} of the needed pairs and the first \m{3{,}000}
unlock \m{58.6\%}.

\subsection{Models: generated reasoning, extracted reasoning, and direct training}

Three training regimes are compared, and the paper's framing depends on keeping them distinct.

\emph{Direct LLM training} is the baseline: one adapter, one sequence, rationale and SMILES emitted
together, with nothing in the objective forcing the molecule to depend on the rationale and no handle
to intervene on under a mismatched-intent test. \emph{Generated reasoning} is our $R$, a separate
adapter emitting the intent from the complex alone and supplying $M$ at inference.
\emph{Extracted reasoning} is $R$'s supervision drawn from the literature rather than from the
agent-written intents of the crystal and potency corpora --- a difference that is not cosmetic, since
an agent writing an intent for a pair has the product molecule in view, whereas a paper's author
wrote the hypothesis without knowing what the reader would later be told. We treat agent-written intents as a \emph{weak upper bound} on supervision
quality and literature intents as the target, and the comparison between them is a result
(\S\ref{sec:agentic}), not a design choice.

\paragraph{Bases and finetuning.} The primary base is ChemDFM-v2.0-14B~\citep{zhao2024chemdfm};
ether0-24B~\citep{narayanan2025ether0}, Llama-3.1-8B-Instruct~\citep{dubey2024llama3},
SmileyLlama-8B~\citep{cavanagh2024smileyllama} and DESRO-8B~\citep{liu2026desro} are comparators.
Because these differ in scale (8B to 24B), cross-base rows are confounded with capacity and the only
clean contrast is each base against \emph{itself}, with and without our adapter. All arms use LoRA
($r{=}16$, $\alpha{=}32$, dropout 0.05, all seven projections; lr $10^{-4}$ cosine, batch 1 $\times$
accumulation 8, bf16) with one shared configuration across every cell of an ablation, and a single
fixed instruction format for every model rather than per-model chat templates, so that template is
never confounded with checkpoint; the blind arm substitutes ``No pocket structure is provided.'' for
the table. Loss is masked to the response --- the prompt runs a median of roughly 960 tokens against
roughly 100 for the answer --- and the length limit acts as a refusal threshold rather than a
truncation, because truncating would delete contact-table rows from precisely the arm the pocket
ablation is meant to favour.

\paragraph{Leakage controls.} For $M$, any SMILES line inside the provided reason is stripped and
the removal is asserted per row; the analogue SMILES, the analogue's PDB identifier and
analogue-only residue tokens are asserted absent from every prompt. Special-token identifiers are
pinned per vocabulary and re-verified at load.

\paragraph{Checkpoint selection, and the resulting caveat.} Training runs to at most 10 epochs with
early stopping on validation loss, patience 2, and the best checkpoint retained. Consequently, and
stated verbatim in every validation caption:

\begin{quote}
Because both the stopping epoch and the retained checkpoint were selected on the validation split,
every validation number reported here is a model-selection score, not a held-out test score, and is
optimistically biased by an unknown amount.
\end{quote}

\subsection{Splits}

Our held-out claims rest on a \textbf{molecule-level} holdout: neither the input nor the reference
analogue appears anywhere in training. On the 21k corpus this leaves \m{1{,}469} rows
(\m{553} on targets never seen in training, \m{916} on trained targets with new molecules), from
which \m{400} prompts per arm are drawn so that all arms answer identical inputs. Unseen-\emph{target}
holdout is infeasible here: target space repeats heavily across papers, so demanding unseen targets
leaves roughly \m{21} of them and the variance would be dominated by which handful of proteins
landed in the holdout. \TBD{14}{split reconciliation}: the corpus builder inherits a target-level
val fraction (32 val targets) while the holdout above is molecule-level; Methods must state one
authoritative sentence for the split the model was \emph{trained} under and the holdout it was
\emph{scored} on. We do not use random pair splits at all: sibling pairs off one complex share the
pocket, the numbering and a near-identical table, so a random split places them on both sides and
measures memorisation --- in a sibling project on the same pair machinery, \m{71\%} of a random
split's test pairs were exact duplicates of training pairs.

\subsection{Evaluation}

\paragraph{Molecular metrics.} \emph{Validity} is a SMILES extracted and parsed, over all rows, so
a model emitting no molecule scores 0 rather than NaN; it is a gate, not a result.
\emph{Uniqueness} is distinct canonical SMILES over valid ones, and functions here as a
collapse detector --- note that in our arms it \emph{anti}-correlates with quality, the mismatched
arm having both the highest uniqueness (\m{0.937}) and the worst molecules. \emph{TC$\rightarrow$input}
(ECFP4, 2048 bits) is a window check, not a quality score: reference analogues sit near \m{0.68},
which is a band and not a maximum, and \m{1.000} means the input was echoed.
\emph{TC$\rightarrow$reference} is the only metric whose ground truth is a real human design
decision, and it is necessary but not sufficient --- it penalises a different \emph{valid} analogue
--- so it is used only in paired contrasts between arms answering the same input.
\emph{Heavy-atom $\Delta$ MAE} catches the right-size/wrong-edit failure that Tanimoto hides.

\paragraph{The groundedness ladder.} ``Groundedness'' as commonly reported means only that a cited
residue exists in that pair's own table. That is rung L0 of a ladder, and the paper reports the
ladder rather than the headline: L0 existence; L1 residue and distance co-occurring in one row;
L2 interaction type matching the row's typing; L3 the covalent gate above; L4 \emph{site adjacency},
the cited residue contacting at least one atom of the changed set; L5 directional ownership, claims
about what is given up mapping to changed atoms; L6 \emph{atom ownership}, the row's atom name
mapped through the chemical-component graph to a specific input index inside the changed set. No
claim in this paper headlines a rung below L4. Each rung is reported as a pair --- pass rate
\emph{and} measurable fraction --- because the mapping fails on some pairs, and an earlier panel
scoring 6/6 against a rigorous 2/6 was partly a coverage gap rather than a disagreement.
Stereochemistry-only pairs are unscorable at L4 and are reported as unscored, never as failed.
\emph{Invented chemistry} --- a sentence asserting a contact absent from that pair's table --- is the
one rationale metric that needs no panel, and it is monotone across our arms.

\paragraph{Residue specificity, not selectivity.} The overlap between the residues a sentence names
and those the reference sentence names is reported as \emph{residue specificity}. We previously
called it selectivity, which invites the reading ``off-target selectivity'' and the demand for a
kinase panel; it is nothing of the kind. It is printed beside its zero-parameter baseline --- name
the nearest table residue --- because an arm at that baseline is no better than a distance sort.

\paragraph{Panels, and what they cannot do.} Sentence quality is scored both by model labellers
under a fixed rubric and by chemists. Panel-to-panel drift exceeds the between-arm differences we
measure, so every panel number is reported with its seeded controls read first, no comparison is
ever made across panels, and panel figures are never pooled with programmatic ones. We report
\emph{pooled multi-target $R^2$ nowhere}: between-target potency offsets inflate it even when
within-target ranking is random. Per-target averaged statistics with $n$ per target only.

\section{Results}

\subsection{Does the stated intent decide the molecule?}
\label{sec:intent}

All arms answer the same \m{400} prompts from the molecule-level holdout, and the only thing that
differs between them is the sentence handed to $M$. Every comparison is paired on the prompt.

\begin{table}[H]\centering\small
\begin{tabular}{llrrrrr}
\toprule
arm & sentence given to $M$ & validity & uniq. & TC$\rightarrow$input & TC$\rightarrow$ref. (med.) & paired $\Delta$ vs mismatch \\
\midrule
reference (dictated) & agent-written, analogue in view & 0.963 & 0.626 & 0.677 & \textbf{0.918} (1.000) & $+0.379$, $p=5\!\times\!10^{-50}$ \\
$R\!\rightarrow\!M$ & $R$'s own generated intent & 0.965 & 0.238 & 0.746 & 0.570 (0.579) & $+0.033$, $p=0.014$ \\
mismatch & another pair's sentence & 0.915 & 0.937 & 0.718 & 0.537 (0.554) & --- (the control) \\
none & nothing & 1.000 & 0.235 & \textbf{1.000} & 0.684 (0.672) & $+0.147$, $p=5\!\times\!10^{-40}$ \\
\bottomrule
\end{tabular}
\caption{Intent conditioning on the molecule-level holdout, $n=400$ prompts per arm, PocketPair-21k.
Heavy-atom $\Delta$ MAE is \m{0.374} (oracle), \m{2.622} ($R\!\rightarrow\!M$) and \m{3.907}
(mismatch): under its own sentence the model moves about one heavy atom, under a foreign one it
drifts almost four. Validation-selection caveat applies.}
\label{tab:arms}
\end{table}

Three readings, in descending order of how much we would stake on them.

\textbf{A correct intent is worth $+0.379$ Tanimoto over a wrong one} (paired,
$p=5\times10^{-50}$), and the reference-intent median of \m{1.000} means that in more than half of
held-out cases the model returns the exact molecule that was made. This arm measures
\emph{transcription fidelity} and not design: the reference sentences were written by an agent with
the analogue in view, and they name a group that actually changed in \m{86.1\%} of rows (correctly
naming \m{2.09} of the \m{2.71} groups that change per pair). There is no oracle in this corpus. \textbf{A wrong intent is worse than
no intent}: the mismatched arm sits \m{0.147} below the no-intent arm ($p=5\times10^{-40}$) and its
validity falls from \m{100\%} to \m{91.5\%}. A model that merely decorated its output with a
rationale would be indifferent to which rationale it was handed; this one is not.

\textbf{The no-intent arm is a diagnostic, not a baseline.} With no sentence the model returns the
input unchanged in \m{100\%} of rows (TC$\rightarrow$input \m{1.000}), so its \m{0.684} is the
input's own similarity to the reference and all of its property deltas are exactly zero. This is an
artefact of pair training rather than a safety property --- trained on pairs with mostly small
edits, $p(y\mid x)$ with no instruction peaks at $y=x$, and the model learned \emph{copy, then
perturb where told}. It is nonetheless informative: zero intent, zero edit, which is direct evidence
that the sentence is the sole driver of structural change. The engineering consequence is that the
system has no sentence-free fallback mode.

\textbf{The reason model, not the molecule model, is the bottleneck.} $R\!\rightarrow\!M$ reaches
\m{0.570}: it clears the mismatched control only by \m{0.033} ($p=0.014$) and retains
$(0.570-0.537)/(0.918-0.537)=\m{8.7\%}$ of the oracle-over-mismatch gap. Its intents are also
repetitive (uniqueness \m{0.238} against oracle's \m{0.626}). In an earlier replication on the
crystal corpus the same ordering held with a smaller spread --- oracle \m{0.698}, $R\!\rightarrow\!M$
\m{0.518}, mismatch \m{0.503}, $+0.180$ at $p=3\times10^{-27}$, and invented-residue rates of
\m{0.0\%} (chemists), \m{13.1\%} ($R$), \m{27.1\%} ($R\!\rightarrow\!M$) and \m{55.7\%} (mismatch).
Two corpora, two splits, one ordering --- and two further points of arithmetic a reader will do
for us if we do not. First, $R\!\rightarrow\!M$ at \m{0.570} sits \emph{below} the no-intent arm at
\m{0.684}. That comparison is not meaningful, because the no-intent arm attains its figure by
returning the input unchanged, but it must be stated rather than left to be discovered. Second, the
two corpora report $R$ differently --- groundedness \m{0.869} and invented chemistry \m{13.1\%} with
residue specificity \m{0.421} on the crystal corpus, against groundedness \m{1.000} and specificity
\m{0.214} on the potency corpus. These are different corpora, different pockets and different
reference sentences, not two readings of one system: the crystal corpus has chemist-written
reference sentences to compare against, which is why its specificity is higher and its groundedness
lower. We report them side by side and never pool them. We therefore present the factorisation as a measured
\emph{upper bound} and as a product claim --- a chemist who can write the hypothesis gets the
structures today --- and not as an end-to-end gain. A direct comparison of the factorised pipeline
against a jointly trained model on one corpus and one split is \TBD{2}{joint vs.\ factorised};
hypothesis: the joint model matches $R\!\rightarrow\!M$ on molecular metrics and loses on every
rationale metric, because nothing in its objective makes the molecule depend on the reason.

\subsubsection{Beyond similarity: does a property-directed intent move that property?}

Tanimoto to a known answer cannot tell a chemist whether a stated goal was achieved. We therefore
classify each row by what its intent \emph{targets}, read from the sentence text rather than from the
assay, and compare the signed change (generated minus input) across arms on exactly those rows.

\begin{table}[H]\centering\small
\begin{tabular}{llrrrrl}
\toprule
the intent targets & property & rows & oracle & $R\!\rightarrow\!M$ & mismatch & verdict \\
\midrule
polarity $\uparrow$ & TPSA & 259 & $+3.62$ & $+4.84$ & $-0.12$ & moves as asked \\
lipophilicity $\downarrow$ & cLogP & 146 & $-0.63$ & $-0.36$ & $-0.16$ & moves as asked \\
size $\downarrow$ & MW & 75 & $-31.25$ & $+0.82$ & $-3.28$ & moves as asked (oracle only) \\
size $\uparrow$ & MW & 88 & $-1.98$ & $-2.66$ & $-10.54$ & flat \\
rigidity $\downarrow$ & rot.\ bonds & 37 & $+0.72$ & $-0.03$ & $+0.06$ & wrong direction \\
basicity $\downarrow$ & basic N & 31 & $-0.03$ & $-0.03$ & $+0.10$ & flat \\
\bottomrule
\end{tabular}
\caption{Property response, 349 paired rows. The no-intent arm is exactly $0.00$ in every cell
because it echoes the input, which is why it is omitted here. Bucket labels are read from the
conditioning sentence and are therefore partly circular with it; the informative comparison is
across arms within a row, not across rows.}
\label{tab:props}
\end{table}

Three of six buckets move in the stated direction and are flat or reversed under a mismatched
intent. Two fail outright. Steerability is real but uneven, and it is strongest exactly where a
single functional-group swap suffices --- adding polarity, shedding lipophilicity, trimming mass ---
and absent where the edit would have to be planned over several atoms. Confidence intervals and the
full bucket counts are \TBD{5}{property CIs}; hypothesis: the two failing buckets stay failing.

\subsection{What the pocket governs}
\label{sec:pocket}

The same corpus and the same held-out prompts, with one difference: whether the per-atom contact
table was present during training and generation.

\begin{table}[H]\centering\small
\begin{tabular}{lrrrr}
\toprule
 & groundedness & residue specificity & validity & TC$\rightarrow$reference \\
\midrule
contact table shown & \textbf{1.000} & \textbf{0.388} & 0.985 & 0.679 \\
contact table withheld & 0.432 & 0.121 & 1.000 & \textbf{0.680} \\
\bottomrule
\end{tabular}
\caption{Pocket ablation, $n=200$ held-out prompts per arm, PocketPair-21k.}
\label{tab:pocket}
\end{table}

The contact table decides \textbf{what the sentence may cite} --- groundedness \m{1.000} against
\m{0.432}, residue specificity \m{0.388} against \m{0.121} --- and it does \textbf{not} decide the
molecule: \m{0.679} with the table, \m{0.680} without it. Two further results point the same way. A
deliberately \emph{wrong} pocket, swapped in at cohort scale, costs \m{0.013} Tanimoto against the
true-analogue noise floor. And adding van der Waals distance annotations to the prompt changed
nothing measurable.

We report this as the central result rather than as a limitation, and we are explicit about what it
does and does not license. It does not license ``pocket-aware generation''. It does license
\emph{grounded attribution}: a target-specific, structurally checkable hypothesis attached to every
proposal, which is what a project team consumes when it has to defend a synthesis list.

Two caveats cut in opposite directions and both belong here. Against us: groundedness is close to
definitional, because the table sits in the prompt, so one arm can copy residue names and the other
must invent them; \m{1.000} versus \m{0.432} partly measures copying from context, and the only
non-trivial figure, residue specificity \m{0.388}, is low. A pocket-\emph{shuffled} control would
still score \m{1.000} by construction and is reported as such. In our favour: whole-molecule
Tanimoto is a poor instrument for this question. A single-edit analogue lives at
TC$\rightarrow$input \m{0.68}--\m{0.75}, where Tanimoto has almost no dynamic range, so
``the pocket did not move the molecule'' is partly a statement about the measurement. The
discriminating experiment is whether the edit lands on the atoms the table flags --- site-of-change
against cited-residue agreement --- and it is \TBD{12}{site-of-change vs.\ cited residue};
hypothesis: agreement exceeds a distance-sorted baseline but by less than the groundedness gap
suggests. Until that number exists, the negative result stands unmitigated.

\paragraph{A mechanism, and a fix.} The omission is specific and it is ours. Fields that describe
\emph{where a group could go} --- open room at an atom, slot shape, exit vectors --- are computed by
our structure layer and then deliberately excluded from the model's input, because they were
computed only at the edit site and therefore leaked it: the exit-vector field alone identified the
edit site in \m{93\%} of rows. The model is consequently handed a description of the current pose and
nothing about available space. A sentence can be grounded in that; a molecule cannot be derived from
it. Computing the same fields at \emph{every} heavy atom carries no site information and is
leakage-safe, and it is the one change we expect to make the pocket touch the molecule
\TBD{15}{room fields at every atom}; hypothesis: TC$\rightarrow$reference is unchanged but
site-of-change agreement improves.

\subsection{Generated molecules with a stated reason, and what finetuning buys}
\label{sec:gen}

\begin{table}[H]\centering\small
\begin{tabular}{llrrrrr}
\toprule
model & training data & validity & uniq. & TC$\rightarrow$input & sensible & sound reasoning \\
\midrule
ChemDFM-14B & none (base) & 0.118 & 0.871 & 0.793 & 10.6\% & 3.9\% \\
ether0-24B & none (base) & 0.225 & 0.847 & 0.884 & 12.7\% & 0.0\% \\
Llama-3.1-8B-Instruct & none (base) & 0.519 & 0.750 & 0.737 & 32.4\% & 6.5\% \\
SmileyLlama-8B & none (base) & 0.882 & 0.753 & 0.582 & 30.6\% & 0.0\% \\
\midrule
ChemDFM-14B & g2 (1{,}485) & 0.893 & 0.752 & 0.658 & 72.4\% & 22.9\% \\
ChemDFM-14B & g3 (2{,}326) & 0.958 & 0.777 & 0.683 & 95.8\% & 29.9\% \\
ChemDFM-14B & g4 (5{,}366) & 0.966 & 0.751 & 0.644 & 89.2\% & 37.6\% \\
ChemDFM-14B & PocketPair-9k (7{,}692) & \textbf{0.973} & 0.757 & 0.658 & 92.4\% & \textbf{46.5\%} \\
ChemDFM-14B & PocketPair-21k & 0.985 & \textbf{0.173} & 0.795 & \TBD{1}{} & \TBD{1}{} \\
\midrule
ether0-24B & PocketPair-9k & 0.855 & 0.768 & 0.639 & 79.6\% & 16.6\% \\
Llama-3.1-8B-Instruct & g3 & 0.771 & 0.743 & 0.594 & 57.8\% & 3.2\% \\
SmileyLlama-8B & g3 & 0.908 & 0.761 & 0.588 & 51.1\% & 3.8\% \\
DESRO-8B & PocketPair-14k & 0.954 & 0.756 & 0.662 & 90.0\% & 33.3\% \\
\bottomrule
\end{tabular}
\caption{Direct joint training. Panel columns come from one model-labeller panel under a fixed
rubric and seed; they are not comparable with panel figures elsewhere in this paper. Bases differ in
scale, so the only clean contrast is each base against itself.}
\label{tab:main}
\end{table}

Two trends and one cost. Validity \textbf{saturates} --- \m{0.893}, \m{0.958}, \m{0.966}, \m{0.973}
across a five-fold increase in data --- while sound reasoning \textbf{keeps climbing} ---
\m{22.9\%}, \m{29.9\%}, \m{37.6\%}, \m{46.5\%} --- and invented chemistry falls from \m{67.7\%} at
base to \m{42.0\%}. The chemistry is a plumbing problem solved early; the reasoning is the thing
data buys, and it had not plateaued when we stopped. The cost is diversity: the potency-trained
model's uniqueness collapses to \m{0.173}, about seventeen distinct ideas per hundred requests, and
$R$'s intents sit at \m{0.238}. On that evidence the system is a renderer of hypotheses, not a source
of them, and we say so rather than reporting uniqueness as a success.

Three numbers in this table will not carry a claim and we flag them here. \m{46.5\%} sound reasoning
and \m{42.0\%} invented chemistry are not ``usable'' by any standard a chemist would accept. The
panel columns rest on model labellers whose drift exceeds these differences. And every row is a
single run; seeded replicates with intervals are \TBD{16}{seed replication}.

\subsection{Agent-written versus literature-extracted reasoning}
\label{sec:agentic}

The intents supervising the crystal and potency corpora were written by an agent with the analogue
in view; the literature intents were written by the authors of the paper. We hold three results
about the former and none yet about the latter. On a cohort round of \m{613} candidates over five
targets, model labellers rated \m{0.61} on ``would make'' but only \m{0.08} on ``explanation
sound'': the molecules were acceptable and the reasons were not. An ownership audit scored \m{2/6}
where a panel scored \m{6/6}. And the field labelled ``specific'' in the agent's intent schema
turned out to measure the template rather than the content ($p=1.000$ on usefulness). Taken
together these say that agent-written reasoning is a weak teacher, which is the motivation for the
literature corpus rather than an aside.

What the literature corpus delivers is \TBD{1}{chemist panel on $R$'s intents} and
\TBD{2}{literature-trained $M$/$R$/$R\!\rightarrow\!M$}; hypotheses: $R$ trained on extracted
intents scores at parity with the agent-trained $R$ on groundedness and above it on usefulness,
while molecular metrics move little, because groundedness is already saturated by the contact table
and usefulness is the dimension the agent labels never had.

\subsection{Retrospective hit-to-lead}
\label{sec:retro}

With a publication-date cutoff and held-out series, the system proposes and justifies analogues of a
known binder and is scored on recovery of the compounds subsequently reported, on agreement with the
authors' own stated intent, and on potency improvement among analogues that were in fact measured,
against both a pocket-blind and a plain enumeration baseline. This experiment is
\TBD{3}{publication-date retrospective} and is gated on the complexes now being folded; it is also
our contamination control, since the base models were pretrained on literature and on ChEMBL that
plausibly include these pairs. Hypothesis: $R\!\rightarrow\!M$ beats the pocket-blind baseline on
intent agreement and lands between enumeration and oracle on recovery.

One caveat carries over from an earlier split study and must be read before any recovery number:
$R$ is the bottleneck by $\m{18.6\times}$, splitting costs roughly $\m{2\times}$ regardless, and the
figures of \m{0.188}/\m{0.703} from that study were \emph{recovery-near}, not exact recovery.

\subsection{A worked covalent programme: ZAP70}
\label{sec:zap70}

On one real covalent programme --- ZAP70, Cys346, \m{280} measured molecules --- a predictor over the
same structural features reaches absolute MAE \m{0.489} and Spearman \m{0.792} with docking into
4K2R plus pretraining, where docking and pretraining contribute additively, against an honestly
rebuilt baseline of MAE \m{0.723}. Earlier versions of the same pipeline reached \m{0.555} and
\m{0.554} (Spearman \m{0.726}). We note two things about this section rather than burying them. It
is a regression model, not $R$/$M$, and therefore evidence about the structural features rather than
about the factorisation. And a previously circulated figure of \m{0.469} for this target came from
one-leg pair leakage worth \m{0.163} MAE: it is withdrawn, it appears nowhere in this paper, and the
honest number is \m{0.723} for the baseline and \m{0.489} for the best model. At \m{280} molecules
every more complex architecture we tried made things worse, which is the most directly actionable
finding we have for programme-scale data: add structure, not parameters.

\section{Related work}

\paragraph{Pocket-conditioned generation.} Structure-based generative models place atoms in a
measured site and several are explicitly interaction-aware: Pocket2Mol~\citep{peng2022pocket2mol},
TargetDiff~\citep{guan2023targetdiff}, DiffSBDD~\citep{schneuing2024diffsbdd},
DrugFlow~\citep{schneuing2025drugflow}, ResGen~\citep{resgen2023}, PocketFlow~\citep{jiang2024pocketflow}
and Lingo3DMol~\citep{feng2024lingo3dmol}, the last predicting non-covalent interaction sites
directly. They do not state a design hypothesis in words, and their evaluation rests on
distributional and docking-derived scores whose adequacy has been
criticised~\citep{cieplinski2023dockwell,harris2023posecheck,buttenschoen2024posebusters}. Recent
work does prompt an LLM with pockets and required residue contacts~\citep{macdougall2026spatial}, so
we do not claim that reasoning LLMs cannot see a pocket; we claim that no prior system reports
whether its stated reasons are checkable against the ligand's own contact table.

\paragraph{Reasoning over molecular optimisation.} Chemistry LLMs and reasoning models now justify
modifications in natural language~\citep{zhao2024chemdfm,narayanan2025ether0,cavanagh2024smileyllama,
fang2023molinstructions,zhang2024chemllm}. Closest to us, DESRO~\citep{liu2026desro} recovers
reasoning steps from outcome data, DrugR~\citep{liu2026drugr} builds rationales backwards from
computed ADMET outcomes, and Logos~\citep{wen2026logos} emits a reason block then a SMILES. All
three generate reason and molecule inside one model and one trace, so the reason cannot be
substituted and the intervention we rely on is unavailable; none conditions on pocket structure.
MolecularCanvas~\citep{dong2026molecularcanvas} separates an LLM plan from an RDKit executor but is
an interactive tool without held-out oracle or mismatch measurements. Earlier work factorised
rationale-to-molecule where the rationale is a \emph{substructure} rather than a
sentence~\citep{jin2020rationale}. Our contribution against this group is the separation into two
independently trained models with a natural-language interface, the structural grounding of that
interface, and the control arms that quantify it.

\paragraph{Instruction-faithfulness controls.} Pocket-shuffled controls are standard in
structure-based design, and UTGDiff~\citep{utgdiff2024} argues that distributional metrics cannot
detect cross-matched instruction/molecule pairs. Our three-arm decomposition adds the finding that a
mismatched instruction falls \emph{below} the no-instruction arm, and that the no-instruction arm is
degenerate because it echoes the input.

\paragraph{Lead optimisation from measured pairs.} Delta and pairwise models exploit analogue pairs
directly~\citep{fralish2023deepdelta,feng2024actfound}, matched molecular series have long been used
to pick the next compound~\citep{wawer2014whatnext}, and REINVENT~4~\citep{loeffler2024reinvent4}
provides the analogue-generation baseline at matched budget. These supply molecules and potency
estimates, not hypotheses.

\paragraph{Literature mining.} LLM extraction from chemistry papers targets reactions, procedures
and entities~\citep{chemminer2024,materials2025ie}. We are not aware of prior work extracting the
stated design rationale of a medicinal-chemistry modification from paper text.

\paragraph{Cofolding as design input.} We take contact tables from cofolded complexes where no
crystal structure exists~\citep{wohlwend2024boltz1,passaro2025boltz2,abramson2024af3,chai2024chai1},
and we adopt the caveats from prospective cofolding
evaluations~\citep{cofold_prospective_2025} rather than treating a predicted pose as a measurement.

\section{Discussion}

\subsection{What we think this shows}

A design intent written in words is a usable interface between a chemist and a generative model. The
evidence is the gap between a correct intent and a wrong one: $+0.379$ Tanimoto to the molecule
chemists actually made, paired, $p=5\times10^{-50}$, with a wrong intent costing both accuracy and
validity relative to no intent at all. That is instruction-following rather than decoration, and it
means the division of labour a project team would actually want --- the chemist states the
hypothesis, the model renders the structures --- is available now, because the oracle arm \emph{is}
that condition.

The pocket's role turned out to be narrower than the premise we started from, and more interesting
than a null. A per-atom contact table determines what a model's explanation may cite and barely
touches what it builds. For the purpose of enumerating analogues this is a negative result. For the
purpose that actually blocks a programme --- choosing five analogues and defending the choice on
Thursday --- an auditable, target-specific hypothesis per proposal is the deliverable, and the table
is what makes it auditable. We think the honest frame is \emph{grounded attribution}, and we think
the mechanism behind the missing structural effect is identified rather than mysterious: the model is
told where the ligand currently sits and nothing about where a group could go.

\subsection{Leakage, and the part of our headline we do not yet own}
\label{sec:leakage}

The strongest number in this paper is also its most fragile, and we would rather decompose it than
defend it. The oracle intents for the crystal and potency corpora were written by an agent that had
the product molecule in view. A sentence written under that condition can name the transformation
itself, in which case $M$ is paraphrasing a specification rather than reasoning from a hypothesis,
and part of \m{0.919} is leakage rather than design. Three facts sharpen the concern. The oracle arm
sits at TC$\rightarrow$input \m{0.677} with heavy-atom $\Delta$ MAE \m{0.374}: roughly one heavy atom
moves, with the relevant residue named in the prompt, which is close to multiple choice with the
answer supplied. The mismatched control is confounded in the same direction, because a foreign
sentence is not merely uninformative but also off-pocket, and its \m{55.7\%} invented-residue rate is
evidence of exactly that. And a median TC$\rightarrow$reference of \m{1.000} means we \emph{reproduce}
compounds that already exist in the literature and plausibly in the patent estate: recovery is not
novelty, and we never claim improvement over a measured compound anywhere in this paper.

The decomposition needs no new training and no new structures: re-run $M$ with (i) oracle sentences
whose chemical content is redacted but whose pocket hypothesis is kept, and (ii) a \emph{same-pocket}
counterfactual intent in place of the foreign-pair sentence, and report the residual
\TBD{17}{leakage decomposition}. Hypothesis: the residual is well under \m{0.919} and well above
mismatch, and that residual --- faithful rendering of a stated intent --- is the claim that survives.
An enumeration control that scores all single-heavy-atom edits of the input by the same sentence, and
reports recovery by edit size, bounds the same thing from the other side
\TBD{11}{enumeration baseline}.

\subsection{Limitations}

Our strongest conditioning result may be partly a leak, as set out above, and the mismatch control is
imperfect for the same reason. The pocket result is narrow: withholding the contact table changes
what the prose may cite (\m{1.000} vs.\ \m{0.432}) but not the molecule (\m{0.679} vs.\ \m{0.680}), a
wrong pocket costs \m{0.013} Tanimoto, distance annotations changed nothing, and a pocket-shuffled
control would score \m{1.000} by construction; we therefore claim grounded attribution and not
pocket-aware generation. The deployable pipeline recovers only \m{8.7\%} of the oracle-over-mismatch
gap and $R$'s intents are repetitive (uniqueness \m{0.238}); the potency-trained model's uniqueness
falls to \m{0.173}. The no-intent arm is not a neutral baseline, since the model echoes its input in
\m{100\%} of those rows. Steerability is uneven and two of six property buckets move the wrong way.

Our literature supervision is SAR pairs, not matched molecular pairs, and \m{60.7\%} of extracted
rationales restate an outcome rather than state an instruction, so part of what we learn is post-hoc
narration --- the same failure we measured at explanation-sound \m{0.08} in agent-written labels.
Stereochemistry is invisible to our site-of-change machinery, so stereo-only pairs are unscored.
Where no crystal structure exists the contact table is a single-sample cofold: a hypothesis about
which residues line the ligand, not a distance measurement, and never a source of covalent geometry.
Most quality judgements come from model labellers whose panel-to-panel drift exceeds the differences
we measure; our chemist panel is pending and will be small; no comparison crosses panels. All results
are single-run unless an interval is given. Our base models were pretrained on literature and ChEMBL
that may contain our test pairs, which the date-cutoff retrospective controls only partially.

Three things a chemist will ask for and will not find here. There is no synthetic-accessibility or
retrosynthesis triage anywhere \TBD{9}{SA and retro triage}, which for a tool proposing compounds to
make is a real gap rather than a refinement. There is no ADMET profile: TPSA, cLogP and molecular
weight are three dials, and hit-to-lead is a multiparameter problem. And there is no novelty check
against prior art for the molecules we generate that are not the reference
\TBD{10}{prior-art novelty}. No molecule in this paper has been synthesised or assayed.

Finally, on our own record: a claim from an earlier system in this programme failed to reproduce
after publication and is treated as unproven, and the ZAP70 figure of \m{0.469} is withdrawn. We
name both because a reader has no other way to know.

\subsection{Reproducibility, ethics and dual use}

We release DOIs, ChEMBL compound keys and text offsets rather than paper text, with quotations at
fair-use length only, and no artefact carrying a publisher watermark; the text-and-data-mining basis
and licensed access are stated. ChEMBL is used under CC-BY-SA with the release version pinned and
derived tables shared alike; all PDB identifiers are listed, and crystal pairs carry no potency. For
agent supervision we release the labeller model, version and date, the extraction brief verbatim, and
the labels themselves, because labellers drift and the recipe alone does not reproduce them. For
model-as-judge we release judge, prompt, temperature and seeds, and state that panel numbers are not
comparable across panels. Splits, LoRA configuration, seeds, GPU-hours and the cofolding cost are
reported. On dual use: this is a generator of bioactive and covalent molecules; we state our output
screening, exclude toxin and select-agent targets, and note that every molecule here is an
unsynthesised, unassayed prediction. Human panel recruitment, compensation, consent and the
IRB-exemption determination are reported with the panel.

\subsection{Next}

In order of expected value: the leakage decomposition and the same-pocket counterfactual, because
they determine what our headline means; the joint-versus-factorised comparison, because it is the
thesis of the paper and is currently untested; room and exit-vector fields computed at every heavy
atom, because that is the one input change we expect to make the pocket touch the molecule; training
$R$ on the instruction subset of the literature corpus; and a blinded chemist panel over
oracle, $R\!\rightarrow\!M$ and mismatch arms with at least two chemists and reported agreement.


\bibliographystyle{unsrtnat}
\bibliography{references}

\appendix
\section{Placeholder register}

Every \texttt{[TBD-$n$]} in the body appears here with the experiment that supplies it and the
result we expect. The register exists so that an experiment which fails costs a placeholder and a
revert, not a rewrite --- and so that no number enters the paper without a stated prior.

\begin{table}[H]\centering\small
\begin{tabular}{p{0.5cm}p{4.6cm}p{4.2cm}p{5.2cm}}
\toprule
id & what is missing & supplied by & hypothesis \\
\midrule
1 & blinded chemist panel on $R$'s intents, and the panel columns of the 21k row &
$\geq$100 items over oracle / $R\!\rightarrow\!M$ / mismatch, $\geq$2 chemists, agreement reported &
$R$ at parity with the paper ceiling on groundedness, below it on usefulness \\
2 & $M$, $R$, $R\!\rightarrow\!M$ trained on the literature triplets; joint vs.\ factorised on one
corpus and split & training once the cofolds land &
literature-trained $R$ improves intent usefulness; molecular metrics move little; the joint model
matches on molecules and loses on rationale \\
3 & publication-date retrospective recovery & literature corpus + cofolds; doubles as the
contamination control & $R\!\rightarrow\!M$ beats a pocket-blind baseline on intent agreement, lands
between enumeration and oracle on recovery \\
4 & intent identifiability: $k=3$ contrasting intents on one input, blind molecule$\rightarrow$intent
matching & automated diagonal-dominance version first & diagonal dominance significant but well
under 1.0 \\
5 & property buckets with counts and intervals & rerun with per-arm CIs & size-up and rigidity stay
wrong-signed \\
6 & contact-table stability under sampling & re-run a random subset at $\geq$5 diffusion samples;
report residue-set Jaccard, median and p90 $|\Delta d|$, and the fraction of interaction-type labels
that flip; calibrate on the 27 crystal-held complexes & residue sets stable, distances not;
conclusions unchanged \\
9 & synthetic accessibility and retrosynthesis triage & SA + retro hit-rate on $\sim$100 generated
molecules, chemist-triaged & a usable majority, with a measurable tail that is not makeable \\
10 & novelty of generated non-reference molecules against prior art & date-cutoff prior-art check &
a substantial fraction are known compounds; recovery $\neq$ novelty \\
11 & enumeration control & score all single-heavy-atom edits of the input by the same sentence;
report set size and recovery by edit size (0--1 / 2--3 / 4+) & oracle recovery concentrates in the
0--1 bucket \\
12 & site-of-change vs.\ cited-residue agreement & MCS changed-set against the cited rows &
agreement above a distance-sorted baseline, below what the groundedness gap suggests \\
13 & training on the instruction subset of the literature corpus & split the corpus by the
instruction/outcome label and train both & the instruction subset trains a more useful $R$ at equal
token count \\
14 & split reconciliation: target-level val fraction in the corpus builder vs.\ the molecule-level
holdout; 32 vs.\ 21 targets & one authoritative pass over the builder and the scorer & no leakage,
but the sentence in Methods must state both \\
15 & room, slot shape and exit vectors computed at \emph{every} heavy atom rather than at the edit
site & recompute the table; leakage-safe by construction & TC$\rightarrow$reference unchanged;
site-of-change agreement improves \\
16 & seeded replicates with intervals on the main table, the pocket ablation and the four intent arms
& $\geq$3 seeds & orderings hold; several single-run differences do not \\
17 & leakage decomposition of the oracle arm & re-run $M$ on (i) oracle sentences with chemical
content redacted and the pocket hypothesis kept, (ii) a same-pocket counterfactual intent in place
of the foreign-pair sentence & the residual sits well under 0.919 and well above mismatch; that
residual is the claim \\
\bottomrule
\end{tabular}
\end{table}

\section{Contributor memos}

The draft was assembled from five independent expert memos, kept in the repository under
\texttt{manuscript/reviews/}: machine learning (formal setup, training, splits, ablation table,
metric definitions), medicinal chemistry (which results change practice, results ordering, figures,
what a chemist will find missing), literature and differentiation (reference list with verification
status, verdict on each novelty claim, closest prior art, sentences we must not write), area chair
(advance meta-review, claim/evidence audit, required experiments, page budget), and structural
chemistry (adequacy of the contact table, cofold claims, covalent typing, the groundedness ladder).
Disagreements that remain open are recorded in the register above rather than resolved in prose.


\end{document}
